% !TEX program = xelatex
% 数学物理方法2 · §8.3 柱坐标分离变量与贝塞尔方程（精讲讲义）
\documentclass[11pt]{ctexart}
%=====================================================================
% theme-violet.tex —— 讲义精讲 · 紫花主题
% 用法：\documentclass[11pt]{ctexart}
%       \input{assets/theme-violet}
%       \begin{document} ... \end{document}
% 只用 ctexart 自带能力 + xcolor/enumitem/fancyhdr/pifont/amsmath，
% 不依赖 tikz/tcolorbox，保证 MiKTeX 可直接编译。
%=====================================================================

\usepackage[a4paper,left=2.0cm,right=2.0cm,top=2.1cm,bottom=2.3cm,%
            headsep=12pt,footskip=22pt]{geometry}
\usepackage{amsmath,amssymb}
\usepackage{needspace}
\usepackage{xcolor}
\usepackage{array}
\usepackage{enumitem}
\usepackage{pifont}
\usepackage{fancyhdr}
\usepackage[hidelinks]{hyperref}

%---------------------------------------------------------------- 配色
\definecolor{violet}{RGB}{107,78,143}      % 主色：紫
\definecolor{violetdeep}{RGB}{63,44,92}    % 深紫：标题/强调
\definecolor{petal}{RGB}{183,159,208}      % 花瓣：装饰线、花饰
\definecolor{bandviolet}{RGB}{247,243,251} % 浅紫底：档案/提示框
\definecolor{bandcream}{RGB}{252,249,242}  % 米白底：易错/注意框
\definecolor{rulepurple}{RGB}{217,207,230} % 分隔线
\definecolor{inkgray}{RGB}{92,96,110}      % 次要文字

%---------------------------------------------------------------- 字体
% 有 Noto CJK 用 Noto，没有就退回 ctex 默认（Windows 下为宋体/黑体）
\IfFontExistsTF{Noto Serif CJK SC}{\setCJKmainfont{Noto Serif CJK SC}}{}
\IfFontExistsTF{Noto Sans CJK SC}{\setCJKsansfont{Noto Sans CJK SC}}{}

%---------------------------------------------------------------- 版式
\setlist{nosep,leftmargin=1.9em}
\linespread{1.28}
\setlength{\parskip}{3.5pt}
\setlength{\parindent}{0pt}
\setlength{\tabcolsep}{5pt}
\renewcommand{\arraystretch}{1.28}

%---------------------------------------------------------------- 花饰
% \ding{95} = 五瓣小花，\ding{96} = 星形花；都来自 pifont，必定存在
\newcommand{\orn}{\textcolor{petal}{\ding{95}}}          % 行内小标
\newcommand{\ornstar}{\textcolor{petal}{\ding{96}}}      % 星形花
\newcommand{\flowerrule}{%                                % 居中花饰分隔线
  \par\vspace{7pt}%
  \noindent\makebox[\linewidth]{\color{petal}%
    \rule{0.20\linewidth}{0.6pt}\hspace{8pt}{\small\ding{95}}\hspace{4pt}%
    \raisebox{1.1pt}{\tiny\ding{95}}\hspace{4pt}{\small\ding{95}}\hspace{8pt}%
    \rule{0.20\linewidth}{0.6pt}}%
  \par\vspace{7pt}}

%---------------------------------------------------------------- 标题与节头
\newcommand{\doctitle}[3]{%   #1 主标题  #2 副标题  #3 元信息（可用 \\ 分行）
  \begin{center}
    {\color{petal}\small \orn\enspace #2\enspace \ornstar}\\[7pt]
    {\LARGE\bfseries\color{violetdeep} #1}\\[10pt]
    {\color{petal}\rule{0.55\linewidth}{1pt}}\\[4pt]
    {\small\color{inkgray} #3}
  \end{center}\vspace{2pt}}

\newcommand{\hdr}[2]{%   #1 章节号（如 §0）  #2 标题
  \par\vspace{12pt}%
  \noindent\colorbox{violetdeep}{\color{white}\bfseries\small\;#1\;}\hspace{7pt}%
  {\color{violetdeep}\large\bfseries #2}\par\nobreak
  \vspace{3pt}{\color{petal}\hrule height 1.2pt}\vspace{6pt}}

\newcommand{\subhdr}[1]{%
  \par\vspace{6pt}%
  \noindent{\color{violet}\bfseries #1}\par\vspace{1pt}}

%---------------------------------------------------------------- 提示框
\newcommand{\callout}[1]{% 浅紫底：讲义档案、读法说明
  \par\medskip\noindent
  \fcolorbox{rulepurple}{bandviolet}{\parbox{\dimexpr\linewidth-2\fboxsep-2\fboxrule}{\vspace{2pt}#1\vspace{2pt}}}%
  \par\medskip}

\newcommand{\warnbox}[1]{% 米白底：易错、边界条件
  \par\medskip\noindent
  \fcolorbox{rulepurple}{bandcream}{\parbox{\dimexpr\linewidth-2\fboxsep-2\fboxrule}{\vspace{2pt}#1\vspace{2pt}}}%
  \par\medskip}

\newcommand{\gapitem}[2]{% #1 缺口类型（紫） #2 内容
  \par\vspace{2pt}\noindent
  \hangindent=1.4em\hangafter=1
  {\color{petal}\small\ding{96}}\enspace{\bfseries\color{violetdeep}#1}\enspace #2\par}

% 题干块（\qstem）：每道例题／作业的「详解」开头必须先复述题干与出处。
% 起因：2026-10-25 第五章讲义把七道作业的题面集中放在一页、详解另起一页，
% 读者翻到「第 4 题」时看不到条件，等于一道没有题干的题目。
% 规则：详解块内自带题干，读者不翻页就能读懂题在问什么。
\newcommand{\qstem}[1]{%
  \par\medskip\noindent
  \fcolorbox{violet}{bandviolet}{%
    \parbox{\dimexpr\linewidth-2\fboxsep-2\fboxrule}{%
      \vspace{2pt}%
      {\small\bfseries\color{violet}题干}\hspace{0.8em}#1%
      \vspace{2pt}}}%
  \par\medskip}

% v2: 保护题干与开头，长题可跨页，续页补上下文。
\newcommand{\qbreak}{\needspace{6\baselineskip}}
% 独立练习：ID、来源、题干、考点、答案标签（通常 ans:ID）。
\newcommand{\practice}[5]{%
  \qbreak\subhdr{练习 #1}\qstem{（#2）#3}%
  \textbf{考点}：#4\par
  详解见附录第~\pageref{#5}~页。\par\medskip}
% 附录开头：ID、来源、条件摘要；自动产生 ans:ID 标签。
\newcommand{\answerstart}[3]{%
  \qbreak\subhdr{#1\quad 详解}\label{ans:#1}\qstem{（#2）#3}}
\newcommand{\answercontinue}[2]{%
  \needspace{5\baselineskip}\subhdr{#1\quad 续解}%
  {\small\color{inkgray}本页接续：#2}\par\medskip}
% 延迟提供默认值，让旧讲义可在导言区保留自己的 \newcommand{\step}。
\AtBeginDocument{\providecommand{\step}[1]{\needspace{3\baselineskip}\par\medskip\textbf{#1}\par}}

% 前提旁注：只放在「需要理解的那一步」后面，给陈述 + 本题代入，不重证。
\newcommand{\aside}[1]{%
  \par\vspace{1pt}\noindent
  \hangindent=1.6em\hangafter=1
  {\small\color{inkgray}%
    \hspace*{0.6em}{\color{petal}$\triangleright$}\hspace{2pt}#1}%
  \par\vspace{2pt}}

%---------------------------------------------------------------- 表格
% 本机 MiKTeX 的 colortbl 与 array 版本不匹配（\rowcolor 会报错），
% 故表格统一用「上粗紫线 + 表头 + 中细线」的三线式，不依赖 colortbl。
\newcommand{\tabtop}{\par\vspace{3pt}{\color{violet}\hrule height 1.2pt}\vspace{1pt}}
\newcommand{\tabmid}{\par\vspace{1pt}{\color{petal}\hrule height 0.5pt}\vspace{1pt}}
\newcommand{\tabbot}{\par\vspace{3pt}{\color{rulepurple}\hrule height 0.5pt}\vspace{8pt}}
\newcommand{\thead}[1]{{\bfseries\color{violetdeep}#1}}

%---------------------------------------------------------------- 页眉页脚
\pagestyle{fancy}
\fancyhf{}
\renewcommand{\headrulewidth}{0.4pt}
\renewcommand{\headrule}{\hbox to\headwidth{\color{rulepurple}\leaders\hrule height \headrulewidth\hfill}}
\fancyhead[L]{\small\color{inkgray}\songti \doctag}
\fancyhead[R]{\small\color{inkgray}\songti \coursetag}
\fancyfoot[C]{\small\color{inkgray}\orn\enspace{\color{violet}\thepage}\enspace\orn}

% 每篇讲义在导言区设这两个值
\newcommand{\doctag}{}
\newcommand{\coursetag}{}
\newcommand{\settags}[2]{\renewcommand{\doctag}{#1}\renewcommand{\coursetag}{#2}}

%---------------------------------------------------------------- 数学记号
\newcommand{\vect}[1]{\boldsymbol{#1}}
\newcommand{\dd}{\,\mathrm{d}}
\newcommand{\ddx}{\frac{\mathrm{d}}{\mathrm{d}x}}
\addtolength{\jot}{2pt}



\usepackage{tikz}
\usetikzlibrary{arrows.meta,calc,positioning}
\settags{\S8.3 · 柱坐标与贝塞尔方程}{数学物理方法2}

% 同一张表内的中隔线
\newcommand{\tabmidrule}{\noalign{\vspace{1pt}{\color{petal}\hrule height 0.5pt}\vspace{1pt}}}

% 圈码交给中文字体
\xeCJKDeclareCharClass{CJK}{"2460 -> "24FF}

% 表头不进第二行
\renewcommand{\thead}[1]{\leavevmode{\bfseries\color{violetdeep}#1}}
% 文本列：左对齐，避免窄列被拉出大片空隙
\newcolumntype{L}[1]{>{\raggedright\arraybackslash}p{#1}}

% 图与图注不分页
\newenvironment{figbox}{\par\begin{center}\begin{minipage}{0.96\linewidth}\centering}{\end{minipage}\end{center}\par}
\newcommand{\figcap}[1]{\par\vspace{4pt}{\small\color{inkgray}#1}\par}

% 详解块起页保护
\renewcommand{\qbreak}{\needspace{0.28\textheight}}

% 标题不落页底
\let\origsubhdr\subhdr
\renewcommand{\subhdr}[1]{\par\needspace{3\baselineskip}\origsubhdr{#1}}
\let\orighdr\hdr
\renewcommand{\hdr}[2]{\par\needspace{6\baselineskip}\orighdr{#1}{#2}}

% 定理与证明
\newcommand{\thm}[2]{\par\medskip\noindent{\color{violetdeep}\bfseries #1}\hspace{0.5em}{\bfseries #2}\par\nobreak\vspace{2pt}}
\newcommand{\pfline}{\par\smallskip\noindent\textbf{证明。}\enspace}

% 公式间距
\setlength{\abovedisplayskip}{7pt plus 2pt minus 2pt}
\setlength{\belowdisplayskip}{7pt plus 2pt minus 2pt}
\setlength{\abovedisplayshortskip}{4pt plus 1pt}
\setlength{\belowdisplayshortskip}{5pt plus 1pt minus 1pt}

\begin{document}

\doctitle{\S8.3\quad 柱坐标下的分离变量与贝塞尔方程}
{从圆域里的波动方程出发，到贝塞尔函数的零点、正交性与柱坐标的三重级数}
{依据课件《Chapter 8 -- 20261116.pdf》15 页（p1--p15）\\2026-11-16\quad 数学物理方法2}

\callout{\textbf{这本讲义怎么用。}课件 15 页分成两段：p1--p6 在圆域（两个变量）里做波动与输运方程，p7--p15 进到圆柱（三个变量），中间 p6 是两道作业。整册的主线只有一条：\textbf{分离变量把一个偏微分方程拆成一串常微分方程，最后一个变系数方程就是贝塞尔方程}。

三件事请先知道：

\textbf{一、}课件把贝塞尔函数直接当成已知（p2 只写“Chapter 13”），但“$k$ 由 $J_m$ 的零点决定”这句话，若不知道零点的存在性与正交性，就只是一句口号。所以本讲补了 \S2：贝塞尔方程的级数解、零点的性质、正交性与模、Fourier--Bessel 展开。这是必要前置，不是选读。

\textbf{二、}课件里 $k_n$ 与 $\mu_n$ 的下标只写了 $n$，实际上它们还依赖阶数 $m$（因为 $x_n^{(m)}$ 是 $m$ 阶贝塞尔的第 $n$ 个零点）。本讲一律写 $k_{mn}=x_n^{(m)}/a$，并在 \S6.2 说明这处记号。

\textbf{三、}课件 p9、p10、p11 用 $\mu$ 的符号把柱坐标问题分成三类，这是本节最该带走的一张表（\S4.3）。判断依据不是“喜欢哪种函数”，而是\textbf{哪一组边界是齐次的}：侧面齐次 $\Rightarrow\mu>0$，上下底齐次 $\Rightarrow\mu<0$，侧面是第二类齐次且轴对称 $\Rightarrow\mu=0$。}

\subhdr{目录（页码自动解析）}
\tabtop
\begin{tabular}{@{}L{0.08\linewidth}L{0.78\linewidth}r@{}}
\thead{节} & \thead{内容} & \thead{页} \\
\tabmidrule
\S0 & 主线图、记号、这一节要解决什么 & \pageref{sec:s0} \\
\S1 & 圆域波动方程：两次分离与零点的出场（p1--p3） & \pageref{sec:s1} \\
\S2 & 贝塞尔方程：级数解、零点、正交性与展开（前置） & \pageref{sec:s2} \\
\S3 & 圆域输运方程：同一个空间谱（p4--p5） & \pageref{sec:s3} \\
\S4 & 柱坐标 Laplace：$\mu$ 的三种符号（p7--p12） & \pageref{sec:s4} \\
\S5 & 柱坐标的波与输运：三重级数（p13--p15） & \pageref{sec:s5} \\
\S6 & 速查、易错与记号说明 & \pageref{sec:s6} \\
\S7 & 课件页码对照表 & \pageref{sec:s7} \\
\S8 & 三句话总结 & \pageref{sec:s8} \\
\S9 & 独立练习（四题，含课件作业） & \pageref{sec:practice} \\
附录 A & 四道练习的详解 & \pageref{sec:app} \\
\end{tabular}
\tabbot

%=====================================================================
\hdr{\S0}{主线图、记号、这一节要解决什么}
\phantomsection\label{sec:s0}
%=====================================================================

\S8.2 处理的是圆域上的 Laplace 方程。那里没有时间，只有一个分离常数 $\lambda$，最后落到欧拉方程，解是 $r^{\pm m}$ 与 $\ln r$——全部是初等函数。本节一上来就不同：圆域里出现\textbf{时间}，第一次分离给出一个常数 $k$，第二次分离后剩下的径向方程变成

\[
\frac{\mathrm d^2R}{\mathrm dx^2}+\frac1x\frac{\mathrm dR}{\mathrm dx}+\Big(1-\frac{m^2}{x^2}\Big)R=0,\qquad x=kr,
\]

它没有初等函数解。这就是贝塞尔函数出场的原因：\textbf{不是我们想引入新函数，而是方程的解只能这样表示}。

\begin{center}
\begin{tikzpicture}[x=1cm,y=1cm,
  ovcore/.style={draw=violetdeep,fill=bandviolet,rounded corners=3pt,align=left,inner xsep=6pt,inner ysep=4pt,line width=1pt,font=\footnotesize},
  ovbox/.style={draw=violet,fill=bandviolet,rounded corners=2.5pt,align=left,inner xsep=5pt,inner ysep=3pt,line width=0.7pt,font=\footnotesize},
  ovcream/.style={draw=violet,fill=bandcream,rounded corners=2.5pt,align=left,inner xsep=5pt,inner ysep=3pt,line width=0.7pt,font=\footnotesize},
  ova/.style={-{Stealth[length=2.2mm,width=1.6mm]},violet,line width=0.9pt},
  ovline/.style={violet,line width=0.9pt}]

\node[ovcore,text width=15.0cm,anchor=north west] (CORE) at (0.2,0.05) {%
\textbf{\color{violetdeep}要解决的事}\quad 圆域、圆柱内的波动（$u_{tt}=a^2\nabla^2u$）、输运（$u_t=a^2\nabla^2u$）与稳态（$\nabla^2u=0$）。\\
\textbf{难点}\quad 边界是 $r=$ 常数；分离之后径向方程变系数，没有初等解。};

\node[ovbox,text width=7.2cm,anchor=north west] (A) at (0.2,-1.85) {%
\textbf{\color{violetdeep}\S1\quad 第一次分离：$u=T(t)V(\mathbf r)$}\\[1pt]
波的：$T''+a^2k^2T=0$，空间的：$\nabla^2V+k^2V=0$（Helmholtz）\\
输运的：$T'+a^2k^2T=0$，空间方程相同\\[1pt]
{\color{inkgray}同一个空间问题，两种时间行为}};

\node[ovbox,text width=7.2cm,anchor=north west] (B) at (7.8,-1.85) {%
\textbf{\color{violetdeep}\S2\quad 前置：贝塞尔方程}\\[1pt]
级数解 $J_m$、第二解 $N_m$、零点 $x_n^{(m)}$\\
带权 $r$ 的正交性与模，Fourier--Bessel 展开\\[1pt]
{\color{inkgray}没有这一段，“$k$ 由零点定”无从谈起}};

\node[ovbox,text width=7.2cm,anchor=north west] (C) at (0.2,-4.65) {%
\textbf{\color{violetdeep}\S1.4\quad 边界定本征值}\\[1pt]
圆内有界 $\Rightarrow$ 丢掉 $N_m$；$r=a$ 齐次 $\Rightarrow J_m(ka)=0$\\
$\Rightarrow k_{mn}=x_n^{(m)}/a$，$n=1,2,\dots$\\[1pt]
{\color{inkgray}谱的离散性来自零点，而不是来自区间长度}};

\node[ovbox,text width=7.2cm,anchor=north west] (D) at (7.8,-4.65) {%
\textbf{\color{violetdeep}\S4\quad 柱坐标：第三次分离}\\[1pt]
$u=R\Phi Z$：$\lambda=m^2$，$Z''-\mu Z=0$，径向方程随 $\mu$ 变\\
$\mu>0$ 贝塞尔 \quad $\mu<0$ 修正贝塞尔 \quad $\mu=0$ 欧拉\\[1pt]
{\color{inkgray}选哪一种，由“哪组边界齐次”决定}};

\node[ovcream,text width=15.0cm,anchor=north west] (JOIN) at (0.2,-7.55) {%
\textbf{\color{violetdeep}最后一步：叠加成三重级数，用正交性定系数}\quad
波/输运的谱是 $k_{mn}$（零点）与 $n\pi/h$（底面）的组合；稳态的谱是 $\mu_j$ 与 $\nu_n$。\\
{\color{inkgray}时间因子决定“振荡还是衰减”，空间因子决定“级数展开用什么基”——两者不要混。}};

\node[ovbox,text width=7.2cm,anchor=north west] (E) at (0.2,-9.75) {%
\textbf{\color{violetdeep}\S3\quad 圆域输运（p4--p5）}\\[1pt]
同一套 $J_m(k_{mn}r)\cos m\varphi$，时间换成 $e^{-a^2k^2t}$\\[1pt]
{\color{inkgray}初值条件给出系数，边界条件给出谱}};

\node[ovbox,text width=7.2cm,anchor=north west] (F) at (7.8,-9.75) {%
\textbf{\color{violetdeep}\S5\quad 柱坐标波/输运（p13--p15）}\\[1pt]
Helmholtz：$\lambda$、$\nu^2$、$\mu=k^2-\nu^2$ 三组本征值\\
$k_{nj}^2=\mu_j+\nu_n^2$\\[1pt]
{\color{inkgray}三重和：$m$（角向）、$n$（轴向）、$j$（径向）}};

\draw[ova] (CORE.south -| A.north) ++(0,0.05) -- (A.north);
\draw[ova] (CORE.south -| B.north) ++(0,0.05) -- (B.north);
\draw[ova] (A.south) -- (C.north);
\draw[ova] (B.south) -- (D.north);
\draw[ovline] (C.south -| JOIN.north) ++(0,0.05) -- (JOIN.north);
\draw[ovline] (D.south -| JOIN.north) ++(0,0.05) -- (JOIN.north);
\draw[ova] (JOIN.south -| E.north) ++(0,0.05) -- (E.north);
\draw[ova] (JOIN.south -| F.north) ++(0,0.05) -- (F.north);
\end{tikzpicture}
\figcap{图 0-1\quad 本节主线：先分离时间，再分离空间；空间部分的角向给 $m$，轴向给 $\nu$，径向给贝塞尔方程；本征值由齐次边界条件（零点或 $n\pi/h$）确定。}
\end{center}

\callout{\textbf{读法。}第一次读只要抓住三件事：第一，$k$ 不是任意的，它是“让 $J_m(ka)=0$ 的那一串数除以 $a$”；第二，柱坐标里 $\mu$ 的符号由\textbf{哪一组边界齐次}决定，不由个人偏好决定；第三，所有级数的系数都用同一套正交性公式定，只是基函数从 $\{\sin,\cos\}$ 换成了 $\{J_m(k_{mn}r)\}$。}

\subhdr{0.1\quad 记号}

\begin{center}
\tabtop
\begin{tabular}{@{}L{0.12\linewidth}L{0.80\linewidth}@{}}
\thead{记号} & \thead{含义} \\
\tabmidrule
$a,\ \rho_0$ & 圆/圆柱的半径（课件 p1--p5 用 $\rho_0$，p7 起用 $a$） \\
$h$ & 圆柱的高（p10 起） \\
$k$ & 第一次分离的常数，单位 $1/$长度；波与输运都用它 \\
$k_{mn}$ & $=x_n^{(m)}/a$，第 $m$ 阶贝塞尔第 $n$ 个零点除以半径 \\
$x_n^{(m)}$ & $J_m$ 的第 $n$ 个正零点（$x_1^{(0)}=2.4048$，$x_2^{(0)}=5.5201$，$\dots$） \\
$\lambda$ & 角向分离常数，$\lambda=m^2$ \\
$\mu$ & 柱坐标里 $z$ 方向的分离常数，可为正、负、零 \\
$\nu$ & $\mu=-\nu^2<0$ 时的正数，$\nu_n=n\pi/h$ \\
$J_m,\ N_m$ & 第一类、第二类柱函数（贝塞尔、诺伊曼） \\
$I_m,\ K_m$ & 第一类、第二类修正贝塞尔函数 \\
$\rho_1,\rho_2$ & 环域的内外半径（作业 7） \\
\end{tabular}
\tabbot
\end{center}

%=====================================================================
\hdr{\S1}{圆域波动方程：两次分离与零点的出场（p1--p3）}
\phantomsection\label{sec:s1}
%=====================================================================

\subhdr{1.1\quad 第一次分离：时间与空间（课件 p1）}

在圆域里解齐次波动方程

\begin{equation}
u_{tt}=a^2\nabla^2u\quad\Longleftrightarrow\quad u_{tt}=a^2\Big(\frac1r\frac{\partial}{\partial r}\Big(r\frac{\partial u}{\partial r}\Big)+\frac1{r^2}u_{\varphi\varphi}\Big).
\label{eq:wave}
\end{equation}

设 $u=T(t)V(r,\varphi)$，代入并除以 $a^2TV$：

\[
\frac{T''}{a^2T}=\frac{\nabla^2V}{V}.
\]

左边只含 $t$，右边只含 $r,\varphi$，所以两边都是常数。记它为 $-k^2$（为什么带负号：波动是振动，时间解必须是振荡的 $\cos$、$\sin$，这要求 $T''/T<0$；若取 $+k^2$，时间解变成指数增长或衰减，与“振动”不符）：

\begin{equation}
T''+a^2k^2T=0,
\qquad
\nabla^2V+k^2V=0 .
\label{eq:helm}
\end{equation}

右边叫 \textbf{Helmholtz 方程}。这一步的意义：时间被彻底分离出去，剩下的 $V$ 只由空间边界决定；而 $k$ 是待定的——它要由边界条件挑出一串离散值，这一点在 \S1.4 才落实。

\subhdr{1.2\quad 第二次分离：Helmholtz 的极坐标形式（课件 p1--p2）}

在极坐标下 \eqref{eq:helm} 是

\begin{equation}
\Big[\frac1r\frac{\partial}{\partial r}\Big(r\frac{\partial}{\partial r}\Big)+\frac1{r^2}\frac{\partial^2}{\partial\varphi^2}+k^2\Big]V=0 .
\label{eq:helmPolar}
\end{equation}

再设 $V=R(r)\Phi(\varphi)$，乘以 $r^2/(R\Phi)$：

\begin{equation}
\frac{r}{R}\frac{\mathrm d}{\mathrm dr}\Big(r\frac{\mathrm dR}{\mathrm dr}\Big)+k^2r^2=-\frac{\Phi''}{\Phi}=\lambda .
\label{eq:sep2}
\end{equation}

左边只含 $r$、右边只含 $\varphi$，故 $\lambda$ 是常数。这一步与 \S8.2 的论证一字不差（那里分离 Laplace 算子，这里分离 Helmholtz 算子，区别只是多了 $k^2r^2$ 一项）。于是得到两个方程：

\begin{equation}
\Phi''+\lambda\Phi=0,
\qquad
r\frac{\mathrm d}{\mathrm dr}\Big(r\frac{\mathrm dR}{\mathrm dr}\Big)+(k^2r^2-\lambda)R=0 .
\label{eq:two}
\end{equation}

由周期条件 $\Phi(\varphi+2\pi)=\Phi(\varphi)$，照 \S8.2 定理 1 的论证，$\lambda=m^2$（$m=0,1,2,\dots$），$\Phi_m=A_m\cos m\varphi+B_m\sin m\varphi$。把 $\lambda=m^2$ 代入径向方程：

\begin{equation}
r\frac{\mathrm d}{\mathrm dr}\Big(r\frac{\mathrm dR}{\mathrm dr}\Big)+(k^2r^2-m^2)R=0 .
\label{eq:radial}
\end{equation}

\subhdr{1.3\quad 换元 $x=kr$：贝塞尔方程（课件 p2）}

令 $x=kr$（$k>0$），记 $R(r)=\tilde R(x)$。由 $\dfrac{\mathrm d}{\mathrm dr}=k\dfrac{\mathrm d}{\mathrm dx}$，有 $r\dfrac{\mathrm d}{\mathrm dr}=x\dfrac{\mathrm d}{\mathrm dx}$，代入 \eqref{eq:radial} 并除以 $k^2$：

\begin{equation}
x\frac{\mathrm d}{\mathrm dx}\Big(x\frac{\mathrm d\tilde R}{\mathrm dx}\Big)+(x^2-m^2)\tilde R=0
\quad\Longleftrightarrow\quad
\frac{\mathrm d^2\tilde R}{\mathrm dx^2}+\frac1x\frac{\mathrm d\tilde R}{\mathrm dx}+\Big(1-\frac{m^2}{x^2}\Big)\tilde R=0 .
\label{eq:bessel}
\end{equation}

右边这一形式叫 \textbf{$m$ 阶贝塞尔方程}。它就是课件 p2 中间那一行。它的两个线性无关解记作 $J_m(x)$ 与 $N_m(x)$（第一类、第二类柱函数），于是

\begin{equation}
R_m(r)=C_mJ_m(kr)+D_mN_m(kr).
\label{eq:Rgen}
\end{equation}

到这里，“为什么会出现贝塞尔函数”已经回答完了：它来自 $x^2R''+xR'+(x^2-m^2)R=0$ 这一条变系数方程，而这条方程是 Helmholtz 方程在极坐标下分离的必然产物，不是人为引入的。

\subhdr{1.4\quad 圆内的两个条件：有界性丢 $N_m$，边界定 $k$（课件 p3）}

\textbf{命题 1.1（$k$ 的取值）}\quad 设圆域 $0<r<\rho_0$ 内 $V$ 有界、$2\pi$ 周期，且 $V(\rho_0,\varphi)=0$（第一类齐次边界）。则

\[
k=k_{mn}=\frac{x_n^{(m)}}{\rho_0}\qquad(m=0,1,2,\dots;\ n=1,2,3,\dots),
\]

其中 $x_n^{(m)}$ 是 $J_m$ 的第 $n$ 个正零点；对应的径向函数是 $J_m(k_{mn}r)$。

\pfline

\textbf{第一步：丢掉 $N_m$。}$x\to0$ 时 $N_0(x)\sim\dfrac2\pi\ln\dfrac x2\to-\infty$，$N_m(x)\sim-\dfrac{(m-1)!}{\pi}\Big(\dfrac2x\Big)^m\to\infty$（$m\ge1$）。若 $D_m\ne0$，解在圆心附近无界，与“圆内有界”矛盾。所以 $D_m=0$，只剩 $J_m(kr)$。

\textbf{第二步：用边界条件。}$V(\rho_0,\varphi)=0$ 对一切 $\varphi$ 成立，而 $R_m(\rho_0)=C_mJ_m(k\rho_0)$，故必须有 $J_m(k\rho_0)=0$（$C_m\ne0$，否则解恒为零）。于是 $k\rho_0$ 必须是 $J_m$ 的零点，即 $k\rho_0=x_n^{(m)}$。$\square$

这个命题里有两件事值得单独记住：

\begin{itemize}
  \item \textbf{谱的离散性来自零点，不来自区间长度。}$J_m$ 有无穷多个正零点（\S2.2 会证明这一点），所以 $k$ 有无穷多个取值，且相邻间隔不等距（$x_1^{(0)}=2.405$，$x_2^{(0)}=5.520$，间隔 $3.115$；$x_2-x_1$ 与 $x_3-x_2=3.134$ 并不相等）。这与 \S8.1 的 $n\pi/a$（等距）形成对照。
  \item \textbf{零点依赖阶数 $m$。}$x_1^{(0)}=2.405$，$x_1^{(1)}=3.832$，$x_1^{(2)}=5.136$。所以严格写法是 $k_{mn}$，不能只写 $k_n$。课件 p3 写成 $k_n=x_n^{(m)}/a$，是把 $m$ 当参数看——本讲保留 $k_{mn}$ 以免混淆。
\end{itemize}

\subhdr{1.5\quad 时间方程与通解（课件 p3）}

$k$ 定下来之后，时间方程 $T''+a^2k^2T=0$ 的通解是

\[
T_{mn}(t)=E_{mn}\cos(a k_{mn}t)+F_{mn}\sin(a k_{mn}t).
\]

注意频率是 $a k_{mn}=a x_n^{(m)}/\rho_0$，量纲为 $1/$时间。把三部分乘起来再叠加：

\begin{equation}
u(r,\varphi,t)=\sum_{m=0}^{\infty}\sum_{n=1}^{\infty}J_m(k_{mn}r)\big(A_m\cos m\varphi+B_m\sin m\varphi\big)\big(E_{mn}\cos(a k_{mn}t)+F_{mn}\sin(a k_{mn}t)\big).
\label{eq:wavegen}
\end{equation}

\eqref{eq:wavegen} 里每个 $(m,n)$ 有四个待定常数（$A_m,B_m$ 与 $E,F$ 的乘积可以合并，实际是每个 $(m,n)$ 四个：$\cos m\varphi\cos\omega t$、$\cos m\varphi\sin\omega t$、$\sin m\varphi\cos\omega t$、$\sin m\varphi\sin\omega t$ 的系数）。它们由两个初始条件 $u(r,\varphi,0)=f$、$u_t(r,\varphi,0)=g$ 定，用的就是 \S2.4 的 Fourier--Bessel 系数公式。

\warnbox{\textbf{一个常见误解。}有人把 \eqref{eq:wavegen} 里的 $J_m(k_{mn}r)$ 当成“径向的解”，以为 $m$ 与 $n$ 可以独立取。其实 $k_{mn}$ 的第二个下标 $n$ 就是径向的编号，$m$ 是角向编号；两者通过零点 $x_n^{(m)}$ 绑在一起，不能拆开。}

%=====================================================================
\hdr{\S2}{贝塞尔方程：级数解、零点、正交性与展开}
\phantomsection\label{sec:s2}
%=====================================================================

课件把这一节留给“Chapter 13”。但 \S1.4 已经用到三件事：$J_m$ 有无穷多零点、$J_m$ 在 $x=0$ 有界而 $N_m$ 无界、以及“用 $J_m(k_{mn}r)$ 展开”是合法的。这三件事不补，后面的级数就只是记号的堆叠。本节只补到“够用”为止：级数解的来源、零点的存在性（陈述）、正交性与模、展开公式，最后给一个可以手算核对的例子。

\subhdr{2.1\quad 级数解：$J_m$ 是怎么来的}

\eqref{eq:bessel} 在 $x=0$ 处是正则奇点（系数 $1/x$ 与 $1-m^2/x^2$ 都只有极点）。按 Frobenius 方法设

\[
\tilde R(x)=x^{\rho}\sum_{j=0}^{\infty}c_jx^j\quad(c_0\ne0),
\]

代入 \eqref{eq:bessel} 并令最低次幂的系数为零，得指标方程

\[
\rho^2-m^2=0\quad\Longrightarrow\quad \rho=\pm m .
\]

取 $\rho=m$（保证 $x=0$ 处有界），逐次比较系数得到 $c_1=0$，且

\[
c_{2k}=-\frac{c_{2k-2}}{(2k)(2k+2m)},\qquad
c_{2k}=\frac{(-1)^k}{2^{2k}k!\,(k+m)!}\,c_0 .
\]

（更准确地说，由 $c_{2k}\big[(2k+m)^2-m^2\big]+c_{2k-2}=0$ 逐项解出；这里只列结果，过程是初等的代数。）取 $c_0=2^{-m}$ 作归一化，得到

\begin{equation}
J_m(x)=\sum_{k=0}^{\infty}\frac{(-1)^k}{k!\,(k+m)!}\Big(\frac x2\Big)^{2k+m}.
\label{eq:Jseries}
\end{equation}

这个级数对一切 $x$ 收敛（相邻项之比的极限为 $0$），所以 $J_m$ 是整函数。前几项是

\[
J_0(x)=1-\frac{x^2}{4}+\frac{x^4}{64}-\cdots,\qquad
J_1(x)=\frac x2-\frac{x^3}{16}+\frac{x^5}{384}-\cdots .
\]

\textbf{第二个解。}指标方程的两个根 $\rho=\pm m$ 相差整数 $2m$，当 $m$ 为整数时第二个根给出的级数与 $J_m$ 只差常数倍，必须另找解。第二解记作 $N_m$（诺伊曼函数），它在 $x=0$ 处无界：

\[
N_0(x)\sim\frac2\pi\ln\frac x2,\qquad
N_m(x)\sim-\frac{(m-1)!}{\pi}\Big(\frac2x\Big)^m\ (m\ge1).
\]

这两条渐近式就是 \S1.4 “圆内丢 $N_m$” 的依据。

\subhdr{2.2\quad 零点的存在性与分布}

\thm{命题 2.1（$J_m$ 的零点）}{对每个整数 $m\ge0$，$J_m(x)$ 在 $x>0$ 上有无穷多个零点 $0<x_1^{(m)}<x_2^{(m)}<\cdots$，它们都是单零点，且 $x_n^{(m)}\to\infty$；相邻零点的间隔趋于 $\pi$。}

这里不给出完整证明（它属于常微分方程的振荡理论），但可以说明为什么合理：$J_m$ 满足的方程在 $x\to\infty$ 时趋近于 $\tilde R''+\tilde R=0$，其解是 $\sin x$、$\cos x$，零点间隔趋于 $\pi$；而 $x_n^{(m)}$ 的数值也确实如此（表 2-1）。

\begin{center}
\tabtop
\begin{tabular}{@{}lcccc@{}}
\thead{阶} & \thead{第 1 个} & \thead{第 2 个} & \thead{第 3 个} & \thead{第 4 个} \\
\tabmidrule
$J_0$ & 2.4048 & 5.5201 & 8.6537 & 11.7915 \\
$J_1$ & 3.8317 & 7.0156 & 10.1735 & 13.3237 \\
$J_2$ & 5.1356 & 8.4172 & 11.6198 & 14.7960 \\
\end{tabular}
\tabbot
\par\vspace{4pt}{\small\color{inkgray}表 2-1\quad $J_0,J_1,J_2$ 的前四个正零点}
\end{center}

\begin{figbox}
\begin{tikzpicture}[x=0.62cm,y=2.1cm,>=Stealth]
  \draw[->,inkgray,line width=0.6pt] (0,0)--(12.6,0) node[right,font=\scriptsize]{$x$};
  \draw[->,inkgray,line width=0.6pt] (0,-0.55)--(0,1.12);
  \draw[densely dotted,inkgray] (0,1)--(12.4,1);
  \draw[densely dotted,inkgray] (0,-0.45)--(12.4,-0.45);
  \node[font=\scriptsize,inkgray,left] at (0,1) {$1$};
  \node[font=\scriptsize,inkgray,left] at (0,-0.45) {$-0.45$};
  \draw[violet,line width=1pt] plot coordinates {(0.0,1.000) (0.25,0.984) (0.5,0.938) (0.75,0.864) (1.0,0.765) (1.25,0.646) (1.5,0.512) (1.75,0.369) (2.0,0.224) (2.25,0.083) (2.5,-0.048) (2.75,-0.164) (3.0,-0.260) (3.25,-0.333) (3.5,-0.380) (3.75,-0.401) (4.0,-0.397) (4.25,-0.369) (4.5,-0.321) (4.75,-0.255) (5.0,-0.178) (5.25,-0.093) (5.5,-0.007) (5.75,0.076) (6.0,0.151) (6.25,0.213) (6.5,0.260) (6.75,0.289) (7.0,0.300) (7.25,0.292) (7.5,0.266) (7.75,0.225) (8.0,0.172) (8.25,0.109) (8.5,0.042) (8.75,-0.026) (9.0,-0.090) (9.25,-0.147) (9.5,-0.194) (9.75,-0.227) (10.0,-0.246) (10.25,-0.249) (10.5,-0.237) (10.75,-0.210) (11.0,-0.171) (11.25,-0.123) (11.5,-0.068) (11.75,-0.010) (12.0,0.048)};
  \draw[violetdeep,line width=0.9pt] plot coordinates {(0.0,0.000) (0.25,0.124) (0.5,0.242) (0.75,0.349) (1.0,0.440) (1.25,0.511) (1.5,0.558) (1.75,0.580) (2.0,0.577) (2.25,0.548) (2.5,0.497) (2.75,0.426) (3.0,0.339) (3.25,0.241) (3.5,0.137) (3.75,0.033) (4.0,-0.066) (4.25,-0.156) (4.5,-0.231) (4.75,-0.289) (5.0,-0.328) (5.25,-0.345) (5.5,-0.341) (5.75,-0.318) (6.0,-0.277) (6.25,-0.221) (6.5,-0.154) (6.75,-0.080) (7.0,-0.005) (7.25,0.069) (7.5,0.135) (7.75,0.192) (8.0,0.235) (8.25,0.262) (8.5,0.273) (8.75,0.267) (9.0,0.245) (9.25,0.209) (9.5,0.161) (9.75,0.105) (10.0,0.043) (10.25,-0.019) (10.5,-0.079) (10.75,-0.132) (11.0,-0.177) (11.25,-0.209) (11.5,-0.228) (11.75,-0.233) (12.0,-0.223)};
  \draw[petal,line width=0.9pt,dash pattern=on 4pt off 2pt] plot coordinates {(0.0,0.000) (0.25,0.008) (0.5,0.031) (0.75,0.067) (1.0,0.115) (1.25,0.171) (1.5,0.232) (1.75,0.294) (2.0,0.353) (2.25,0.405) (2.5,0.446) (2.75,0.474) (3.0,0.486) (3.25,0.481) (3.5,0.459) (3.75,0.419) (4.0,0.364) (4.25,0.296) (4.5,0.218) (4.75,0.133) (5.0,0.047) (5.25,-0.038) (5.5,-0.117) (5.75,-0.187) (6.0,-0.243) (6.25,-0.284) (6.5,-0.307) (6.75,-0.313) (7.0,-0.301) (7.25,-0.273) (7.5,-0.230) (7.75,-0.176) (8.0,-0.113) (8.25,-0.046) (8.5,0.022) (8.75,0.087) (9.0,0.145) (9.25,0.193) (9.5,0.228) (9.75,0.249) (10.0,0.255) (10.25,0.245) (10.5,0.222) (10.75,0.185) (11.0,0.139) (11.25,0.085) (11.5,0.028) (11.75,-0.030) (12.0,-0.085)};
  \foreach \x/\lab in {2.4048/{$x_1^{(0)}$},5.5201/{$x_2^{(0)}$},8.6537/{$x_3^{(0)}$}} {
    \fill[violet] (\x,0) circle (1.3pt);
    \node[font=\scriptsize,violet] at (\x,-0.50) {\lab};
  }
  \node[font=\scriptsize,inkgray] at (6.3,-0.95) {$J_0$（实线）\quad $J_1$（深紫）\quad $J_2$（虚线，petal）};
\end{tikzpicture}
\figcap{图 2-1\quad $J_0,J_1,J_2$ 在 $[0,12]$ 上的曲线（按 \eqref{eq:Jseries} 数值计算）。标出的是 $J_0$ 的前三个正零点，它们就是 $k\rho_0$ 的取值。注意零点间隔并不相等。}
\end{figbox}

\subhdr{2.3\quad 正交性与模}

\thm{定理 2.2（贝塞尔函数的正交性）}{设 $k_{mn}=x_n^{(m)}/a$，其中 $x_n^{(m)}$ 是 $J_m$ 的第 $n$ 个正零点。则}

\[
\int_0^aJ_m(k_{ml}r)J_m(k_{mn}r)\,r\dd r=0\quad(l\ne n),
\qquad
\int_0^a\big[J_m(k_{mn}r)\big]^2r\dd r=\frac{a^2}{2}\big[J_{m+1}(x_n^{(m)})\big]^2 .
\]

\pfline

把 \eqref{eq:bessel} 写成自伴形式。令 $Z_n(r)=J_m(k_{mn}r)$，它满足

\[
\frac{\mathrm d}{\mathrm dr}\Big(r\frac{\mathrm dZ_n}{\mathrm dr}\Big)+\Big(k_{mn}^2r-\frac{m^2}{r}\Big)Z_n=0 .
\]

对 $l$ 写同样一式，乘 $Z_n$ 与 $Z_l$ 后相减，得

\[
\frac{\mathrm d}{\mathrm dr}\Big[r\Big(Z_lZ_n'-Z_nZ_l'\Big)\Big]=\big(k_{ml}^2-k_{mn}^2\big)rZ_lZ_n .
\]

两边从 $0$ 到 $a$ 积分。左端等于 $\big[r(Z_lZ_n'-Z_nZ_l')\big]_0^a$：$r=0$ 处因 $Z_n(0)$ 有界、$rZ_n'\to0$ 而为零；$r=a$ 处因 $Z_n(a)=Z_l(a)=0$ 而为零。所以

\[
\big(k_{ml}^2-k_{mn}^2\big)\int_0^aZ_lZ_n\,r\dd r=0 .
\]

当 $l\ne n$ 时 $k_{ml}\ne k_{mn}$（零点互不相同），故积分为零。$\square$

模的公式由 $l\to n$ 的极限得到（或用 $J_m$ 的递推式把 $\partial_kZ$ 化成 $J_{m+1}$），这里直接引用结果；它的用处只有一个：给展开系数配分母。

\subhdr{2.4\quad Fourier--Bessel 展开}

由定理 2.2，$\{J_m(k_{mn}r)\}_{n\ge1}$ 在带权 $r$ 的意义下正交。若 $f(r)$ 在 $(0,a)$ 上分段光滑，则可展开为

\begin{equation}
f(r)=\sum_{n=1}^{\infty}A_nJ_m(k_{mn}r),\qquad
A_n=\frac{2}{a^2\big[J_{m+1}(x_n^{(m)})\big]^2}\int_0^af(r)J_m(k_{mn}r)\,r\dd r .
\label{eq:fb}
\end{equation}

这一步就是“用贝塞尔函数当基”的全部内容。它与 Fourier 级数的关系：若把 $r$ 换成 $\varphi$、$J_m(k_{mn}r)$ 换成 $\cos n\varphi$、权 $r\dd r$ 换成 $\dd\varphi$，公式 \eqref{eq:fb} 在结构上一字不差。

\subhdr{2.5\quad 一个可以手算的例子：$f(r)=1$ 的 $J_0$ 展开}

取 $m=0$、$f(r)=1$。先算积分。由 $\dfrac{\mathrm d}{\mathrm dx}\big[xJ_1(x)\big]=xJ_0(x)$（这可由 \eqref{eq:Jseries} 直接验证），

\[
\int_0^aJ_0(k_{0n}r)\,r\dd r
=\frac1{k_{0n}^2}\int_0^{x_n^{(0)}}J_0(t)\,t\dd t
=\frac1{k_{0n}^2}\Big[tJ_1(t)\Big]_0^{x_n^{(0)}}
=\frac{a^2J_1(x_n^{(0)})}{x_n^{(0)}} .
\]

代进 \eqref{eq:fb}：

\[
A_n=\frac{2}{a^2J_1^2(x_n^{(0)})}\cdot\frac{a^2J_1(x_n^{(0)})}{x_n^{(0)}}=\frac{2}{x_n^{(0)}J_1(x_n^{(0)})}.
\]

前三个系数（用表 2-1 与 $J_1(2.4048)=0.5192$、$J_1(5.5201)=-0.3403$、$J_1(8.6537)=0.2715$）：

\begin{gather*}
A_1=\frac{2}{2.4048\times0.5192}=1.602,\qquad
A_2=\frac{2}{5.5201\times(-0.3403)}=-1.065,\\
A_3=\frac{2}{8.6537\times0.2715}=0.851 .
\end{gather*}

于是

\[
1=\sum_{n=1}^{\infty}\frac{2}{x_n^{(0)}J_1(x_n^{(0)})}\,J_0\Big(\frac{x_n^{(0)}r}{a}\Big),\qquad 0<r<a .
\]

\textbf{两条核对}（做作业时可用来检查自己的系数）：在 $r=0$ 处，$J_0(0)=1$，右边成为 $\sum A_n$，它应收敛到 $1$（部分和 $1.602-1.065+0.851-\cdots$ 在 $1$ 附近摆动，收敛很慢，因为系数只按 $1/\sqrt{x_n}\sim1/\sqrt n$ 衰减）；在 $r\to a^-$ 处，每一项 $J_0(x_n^{(0)})=0$，级数趋于 $0$，而 $f(a^-)=1$——这是端点处必然的跳跃（Fourier--Bessel 级数在端点收敛到左右极限的平均，$(1+0)/2=1/2$），不是算错。

%=====================================================================
\hdr{\S3}{圆域输运方程：同一个空间谱（p4--p5）}
\phantomsection\label{sec:s3}
%=====================================================================

\subhdr{3.1\quad 只换时间方程}

齐次输运方程 $u_t=a^2\nabla^2u$ 做同样的第一次分离 $u=T(t)V(r,\varphi)$：

\[
\frac{T'}{a^2T}=\frac{\nabla^2V}{V}=-k^2 ,
\]

得到

\begin{equation}
T'+a^2k^2T=0,
\qquad
\nabla^2V+k^2V=0 .
\label{eq:transport}
\end{equation}

\textbf{空间方程与波动方程完全相同}（都是 Helmholtz \eqref{eq:helm}），所以 \S1.2--\S1.4 的结论照搬：$V=R_m(r)\Phi_m(\varphi)$、$k=k_{mn}=x_n^{(m)}/\rho_0$、$R_m=J_m(k_{mn}r)$。变的只有时间：

\[
T_{mn}(t)=E_{mn}e^{-a^2k_{mn}^2t}.
\]

于是

\begin{equation}
u(r,\varphi,t)=\sum_{m=0}^{\infty}\sum_{n=1}^{\infty}J_m(k_{mn}r)\big(A_m\cos m\varphi+B_m\sin m\varphi\big)E_{mn}e^{-a^2k_{mn}^2t}.
\label{eq:heatgen}
\end{equation}

系数 $A_mB_mE_{mn}$ 由初始条件 $u(r,\varphi,0)=f(r,\varphi)$ 定：把 $t=0$ 代入，用 \eqref{eq:fb} 逐阶展开。

\begin{center}
\tabtop
\begin{tabular}{@{}L{0.16\linewidth}L{0.34\linewidth}L{0.40\linewidth}@{}}
\thead{环节} & \thead{波动方程} & \thead{输运方程} \\
\tabmidrule
时间方程 & $T''+a^2k^2T=0$ & $T'+a^2k^2T=0$ \\
时间解 & $\cos(a kt)$、$\sin(a kt)$ & $e^{-a^2k^2t}$ \\
初始条件个数 & 两个（$u$ 与 $u_t$） & 一个（$u$） \\
空间部分 & \multicolumn{2}{c}{完全相同：$\nabla^2V+k^2V=0$，$k=k_{mn}$，$V=J_m(k_{mn}r)\Phi_m$} \\
\end{tabular}
\tabbot
\end{center}

\textbf{为什么高阶项衰减得快。}$e^{-a^2k^2t}$ 的指数里是 $k^2$，而 $k_{mn}=x_n^{(m)}/\rho_0$ 随 $n$ 增大。所以 $t$ 稍大以后，只剩 $n=1$（以及 $m=0$ 的那一项）在起作用：温度场迅速变成“最低阶模态”的形状。这是热传导与振动最直观的差别——振动的高阶模态只是频率高，不会自动消失；热传导的高阶模态衰减得极快。

%=====================================================================
\hdr{\S4}{柱坐标 Laplace：$\mu$ 的三种符号（p7--p12）}
\phantomsection\label{sec:s4}
%=====================================================================

\subhdr{4.1\quad 三变量分离（课件 p7）}

柱坐标下的 Laplace 方程是

\begin{equation}
\nabla^2u=\frac1r\frac{\partial}{\partial r}\Big(r\frac{\partial u}{\partial r}\Big)+\frac1{r^2}\frac{\partial^2u}{\partial\varphi^2}+\frac{\partial^2u}{\partial z^2}=0 .
\label{eq:lapcyl}
\end{equation}

设 $u=R(r)\Phi(\varphi)Z(z)$，代入并除以 $R\Phi Z$：

\[
\frac1{rR}\frac{\mathrm d}{\mathrm dr}\Big(r\frac{\mathrm dR}{\mathrm dr}\Big)+\frac1{r^2\Phi}\Phi''+\frac1Z Z''=0 .
\]

乘以 $r^2$：$\dfrac{r}{R}(rR')'+\dfrac{\Phi''}{\Phi}+\dfrac{r^2Z''}{Z}=0$。前两项只含 $r,\varphi$，第三项只含 $z$，于是 $\dfrac{Z''}{Z}=\mu$ 是常数；剩下

\[
\frac{r}{R}(rR')'+\frac{\Phi''}{\Phi}+\mu r^2=0
\quad\Longrightarrow\quad
\frac{\Phi''}{\Phi}=-\lambda,\qquad
\frac1r(rR')'+\Big(\mu-\frac{\lambda}{r^2}\Big)R=0 .
\]

合并成课件 p7 右下那个方程组：

\begin{equation}
\Phi''+\lambda\Phi=0,
\qquad
Z''-\mu Z=0,
\qquad
\frac1r\frac{\mathrm d}{\mathrm dr}\Big(r\frac{\mathrm dR}{\mathrm dr}\Big)+\Big(\mu-\frac{\lambda}{r^2}\Big)R=0 .
\label{eq:cylsys}
\end{equation}

由周期条件，$\lambda=m^2$（$m=0,1,2,\dots$），$\Phi_m=A_m\cos m\varphi+B_m\sin m\varphi$。至此还剩两个常数 $m$ 与 $\mu$，以及径向方程。

\subhdr{4.2\quad 为什么 $\mu$ 的符号决定一切}

把 \eqref{eq:cylsys} 的径向方程写成 $\dfrac1r(rR')'+\Big(\mu-\dfrac{m^2}{r^2}\Big)R=0$。它的“身份”完全由 $\mu$ 的符号决定：

\[
\mu>0:\ \text{令 }x=\sqrt\mu\,r\ \Rightarrow\ \text{贝塞尔方程};
\]
\[
\mu<0:\ \text{令 }x=\sqrt{-\mu}\,r\ \Rightarrow\ \text{修正贝塞尔方程};
\qquad
\mu=0:\ \text{欧拉方程}.
\]

而 $\mu$ 的符号不是自由选的：它由 $Z$ 的方程 $Z''-\mu Z=0$ 与\textbf{上下底面的边界条件}共同决定。$\mu>0$ 时 $Z$ 是双曲函数（指数），$\mu<0$ 时 $Z$ 是三角函数（振荡），$\mu=0$ 时 $Z$ 是线性的。于是有一个非常好用的判据：

\begin{quote}
\textbf{若上下底面给的是齐次条件（$Z$ 必须在这两个 $z$ 值处为零），则 $Z$ 必须振荡，$\mu<0$；若侧面给的是齐次条件（$R$ 必须在 $r=a$ 处为零），则 $R$ 必须振荡，$\mu>0$。}
\end{quote}

这个判据的道理是：\textbf{齐次条件要求对应的函数“回到零”，只有振荡型函数能反复回到零}；指数型函数最多回到零一次（或者根本回不去），给不出无穷多个本征值。

\subhdr{4.3\quad 三种情形的完整解（课件 p9--p11）}

\begin{center}
\tabtop
\begin{tabular}{@{}L{0.11\linewidth}L{0.24\linewidth}L{0.25\linewidth}L{0.29\linewidth}@{}}
\thead{项目} & \thead{$\mu>0$（p9）} & \thead{$\mu=-\nu^2<0$（p10）} & \thead{$\mu=0$（p11）} \\
\tabmidrule
适用 & 侧面 $r=a$ 齐次 & 上下底 $z=0,h$ 齐次 & 侧面为第二类齐次（$\partial_ru=0$）且轴对称 \\
$Z$ 的方程 & $Z''-\mu Z=0$ & $Z''+\nu^2Z=0$ & $Z''=0$ \\
$Z$ 的解 & $E e^{\sqrt\mu z}+F e^{-\sqrt\mu z}$ & $E\sin\nu z+F\cos\nu z$ & $A+Bz$ \\
$Z$ 的本征值 & 由侧面条件定（不需要） & $\nu_n=n\pi/h$ & 无 \\
$R$ 的方程 & 贝塞尔方程 & 修正贝塞尔方程 & 欧拉方程 \\
$R$ 的解 & $CJ_m(\sqrt\mu r)+DN_m(\sqrt\mu r)$ & $CI_m(\nu r)+DK_m(\nu r)$ & $m=0$：$C_1+C_2\ln r$；$m>0$：$C_1r^m+C_2r^{-m}$ \\
$R$ 的本征值 & $\mu_j=\big(x_j^{(m)}/a\big)^2$ & 无（$\nu$ 已由底面定） & 无 \\
\end{tabular}
\tabbot
\par\vspace{4pt}{\small\color{inkgray}表 4-1\quad $\mu$ 的三种符号与被选中的径向方程}
\end{center}

三种情形各有一句“为什么”：

\begin{itemize}
  \item \textbf{$\mu>0$（p9）}：$Z$ 是实指数，不能反复回零，所以 $z$ 方向不给本征值；本征值只能来自侧面的 $J_m(\sqrt\mu\,a)=0$，即 $\sqrt\mu=x_j^{(m)}/a$。注意这里 $Z$ 在 $z$ 方向没有齐次条件，$E,F$ 由上下底的具体（非齐次）数据定。
  \item \textbf{$\mu<0$（p10）}：$Z=\sin\nu z$、$\cos\nu z$ 可以满足 $Z(0)=Z(h)=0$（取 $\sin$，$\nu_n=n\pi/h$）。此时径向方程变成修正贝塞尔方程，$I_m$ 随 $r$ 增长、$K_m$ 随 $r$ 衰减且 $K_m(0)=\infty$；在实心圆柱里必须丢掉 $K_m$。
  \item \textbf{$\mu=0$（p11）}：$Z=A+Bz$，$\Phi$ 最好是常数（轴对称，$\lambda=0$）。径向方程是欧拉方程，$m=0$ 时 $R=C_1+C_2\ln r$；它的导数 $R'=C_2/r$，在 $r=a$ 上等于 $C_2/a$。若侧面条件是第二类齐次 $\partial_ru|_{r=a}=0$，则 $C_2=0$，$R=C_1$ 自动满足——这就是课件 p11 说 “This satisfies the second kind homogeneous boundary condition of the side of the cylinder” 的意思。
\end{itemize}

\begin{figbox}
\begin{tikzpicture}[x=1cm,y=1cm,>=Stealth]
  \draw[violet,line width=1pt] (-1.6,-1.2) -- (-1.6,1.6) arc (180:360:1.6 and 0.5) -- (1.6,-1.2) arc (0:-180:1.6 and 0.5);
  \draw[violet,line width=1pt,dashed] (-1.6,1.6) arc (180:0:1.6 and 0.5);
  \draw[violet,line width=1pt,dashed] (-1.6,-1.2) arc (180:0:1.6 and 0.5);
  \draw[inkgray,densely dotted] (0,-1.2)--(0,1.6);
  \node[font=\scriptsize,inkgray,right] at (0.05,0.2) {$z$};
  \draw[->,violetdeep,line width=0.7pt] (0,-1.2)--(1.25,0.14);
  \node[font=\scriptsize,violetdeep,right] at (0.85,-0.28) {$r$};
  \draw[->,violetdeep,line width=0.5pt] (-2.15,1.05)--(-1.62,0.55);
  \node[font=\scriptsize,violetdeep,anchor=east] at (-2.15,1.05) {$r=a$};
  \node[font=\scriptsize,violetdeep] at (1.95,-1.30) {$z=0$};
  \node[font=\scriptsize,violetdeep] at (1.95,1.92) {$z=h$};
  \node[font=\footnotesize,inkgray,align=left,anchor=west] at (2.15,1.05) {侧面 $r=a$：第一类 $\Rightarrow\mu>0$};
  \node[font=\footnotesize,inkgray,align=left,anchor=west] at (2.15,0.35) {上下底 $z=0,h$：齐次 $\Rightarrow\mu<0$};
  \node[font=\footnotesize,inkgray,align=left,anchor=west] at (2.15,-0.35) {侧面第二类且轴对称 $\Rightarrow\mu=0$};
  \node[font=\footnotesize,inkgray,align=left,anchor=west] at (2.15,-1.5) {本征值：$\mu_j=\big(x_j^{(m)}/a\big)^2$ 或 $\nu_n=n\pi/h$};
\end{tikzpicture}
\figcap{图 4-1\quad 圆柱的三种边界配置。哪一组边界齐次，就由哪一组定本征值，也就决定了 $\mu$ 的符号（表 4-1）。}
\end{figbox}

\subhdr{4.4\quad 通解与“$R_m$、$Z_n$ 随边界条件而变”（课件 p12）}

课件 p12 给出柱坐标 Laplace 的通解

\begin{equation}
u(r,\varphi,z)=\sum_{m=0}^{\infty}\sum_{n=1}^{\infty}R_m(r)Z_n(z)\Phi_m(\varphi),
\label{eq:cylgen}
\end{equation}

并加了一句 “$R_m(r)$ and $Z_n(z)$ are different for different boundary conditions”。这句话是本节的总结，展开就是表 4-1 的三行。具体做题时按四步走：

\begin{enumerate}[label=(\arabic*)]
  \item 看哪一组边界齐次：侧面齐次 $\to$ 用 $\mu>0$ 这一行；上下底齐次 $\to$ 用 $\mu<0$ 这一行；侧面第二类且轴对称 $\to$ 用 $\mu=0$ 这一行。
  \item 写出该行的 $Z(z)$ 与本征值（若这一行没有本征值，就保留常数）。
  \item 写出该行的 $R(r)$，并用“$r\to0$ 有界”或“$r\to\infty$ 有界”删去第二个解（实心柱删 $N_m$ 或 $K_m$）。
  \item 把剩下的级数代入非齐次的那一组边界条件，用正交性逐项定系数。
\end{enumerate}

\warnbox{\textbf{一个容易错的点。}$\mu>0$ 时 $Z$ 用 $e^{\pm\sqrt\mu z}$，但很多题目里上下底是 $u=0$（齐次）。这时不能用 $\mu>0$——因为 $e^{\pm\sqrt\mu z}$ 在 $z=0$ 与 $z=h$ 同时为零只有平凡解。判断的依据始终是“齐次条件要求函数反复回零”，不是“书上画的是哪种”。}

%=====================================================================
\hdr{\S5}{柱坐标的波与输运：三重级数（p13--p15）}
\phantomsection\label{sec:s5}
%=====================================================================

\subhdr{5.1\quad Helmholtz 在柱坐标（课件 p13）}

波动方程 $u_{tt}=a^2\nabla^2u$ 的第一次分离与 \S1.1 相同，给出

\[
T''+a^2k^2T=0,
\qquad
\nabla^2V+k^2V=0 .
\]

把 Helmholtz 方程写在柱坐标（把 \eqref{eq:lapcyl} 的 $0$ 换成 $-k^2V$）：

\begin{equation}
\frac1r\frac{\partial}{\partial r}\Big(r\frac{\partial V}{\partial r}\Big)+\frac1{r^2}\frac{\partial^2V}{\partial\varphi^2}+\frac{\partial^2V}{\partial z^2}+k^2V=0 .
\label{eq:helmcyl}
\end{equation}

设 $V=R\Phi Z$，与 \S4.1 完全同样的推导给出

\begin{equation}
\Phi''+\lambda\Phi=0,\qquad
Z''+\nu^2Z=0,\qquad
\frac1r\frac{\mathrm d}{\mathrm dr}\Big(r\frac{\mathrm dR}{\mathrm dr}\Big)+\Big(k^2-\nu^2-\frac{\lambda}{r^2}\Big)R=0 .
\label{eq:helmsys}
\end{equation}

注意与 \eqref{eq:cylsys} 的差别：$Z$ 的方程变成 $Z''+\nu^2Z=0$（因为 Helmholtz 里 $k^2$ 已经占了“正项”的位置），而径向方程里出现的是组合

\begin{equation}
\mu=k^2-\nu^2 .
\label{eq:mudef}
\end{equation}

只有 $\mu>0$ 时径向方程才是贝塞尔方程；课件 p13 末尾写的就是 $\mu=k^2-\nu^2>0$。

\subhdr{5.2\quad 三组本征值（课件 p14）}

\begin{center}
\tabtop
\begin{tabular}{@{}L{0.18\linewidth}L{0.34\linewidth}L{0.38\linewidth}@{}}
\thead{方向} & \thead{条件} & \thead{本征值} \\
\tabmidrule
角向 & 周期 & $\lambda=m^2$，$m=0,1,2,\dots$；$\Phi_m=\cos m\varphi,\ \sin m\varphi$ \\
轴向 & $Z(0)=Z(h)=0$（底面齐次） & $\nu_n=\dfrac{n\pi}{h}$，$n=1,2,\dots$；$Z_n=\sin\dfrac{n\pi z}{h}$ \\
径向 & $R(a)=0$（侧面齐次） & $\mu_j=\Big(\dfrac{x_j^{(m)}}{a}\Big)^2$，$j=1,2,\dots$；$R_m=J_m(\sqrt{\mu_j}\,r)$ \\
频率 & $k^2=\mu+\nu^2$ & $k_{nj}=\sqrt{\mu_j+\nu_n^2}$；$T_{nj}=G\cos(k_{nj}at)+H\sin(k_{nj}at)$ \\
\end{tabular}
\tabbot
\end{center}

课件 p14 里写 “$k=\sqrt{\mu_j+\nu_n^2}$”，与这里的 $k_{nj}$ 是同一个东西；注意它同时依赖 $j$ 与 $n$（还依赖 $m$，因为 $\mu_j$ 里的零点是 $m$ 阶的）。

\subhdr{5.3\quad 通解（课件 p15）}

把三部分乘起来、对三个指标求和：

\begin{equation}
\begin{aligned}
u(r,\varphi,z,t)=\sum_{m=0}^{\infty}\sum_{n=1}^{\infty}\sum_{j=1}^{\infty}&J_m(\sqrt{\mu_j}\,r)\sin\frac{n\pi z}{h}\\
&\times\big(A_m\cos m\varphi+B_m\sin m\varphi\big)\big(G_{nj}\cos(k_{nj}at)+H_{nj}\sin(k_{nj}at)\big).
\end{aligned}
\label{eq:triple}
\end{equation}

输运方程只是把时间因子换成 $e^{-a^2k_{nj}^2t}$：

\begin{equation}
u(r,\varphi,z,t)=\sum_{m,n,j}J_m(\sqrt{\mu_j}\,r)\sin\frac{n\pi z}{h}\big(A_m\cos m\varphi+B_m\sin m\varphi\big)E_{nj}e^{-a^2k_{nj}^2t}.
\label{eq:tripleheat}
\end{equation}

系数由初始条件（波动还要 $u_t$）与可能存在的非齐次边界值确定，用的仍是 \eqref{eq:fb} 的那套正交性——只是要在 $r$、$\varphi$、$z$ 三个方向各用一次。

\warnbox{\textbf{三重和不是“随便乘”。}三个指标的来源不同：$m$ 来自 $\varphi$ 的周期条件，$n$ 来自 $z$ 的两端条件，$j$ 来自 $r=a$ 的条件。少一个条件，就少一个求和指标。若侧面不是齐次（比如 $u(a,\varphi,z)$ 给定），径向就没有本征值，$\mu$ 要由 $z$ 方向定，级数结构随之改变——这正是表 4-1 存在的理由。}

%=====================================================================
\hdr{\S6}{速查、易错与记号说明}
\phantomsection\label{sec:s6}
%=====================================================================

\subhdr{6.1\quad 六个式子}

\begin{center}
\tabtop
\begin{tabular}{@{}L{0.22\linewidth}L{0.72\linewidth}@{}}
\thead{对象} & \thead{式子} \\
\tabmidrule
柱坐标 Laplace 算子 & $\nabla^2u=\frac1r(ru_r)_r+\frac1{r^2}u_{\varphi\varphi}+u_{zz}$ \\
贝塞尔方程 & $x^2R''+xR'+(x^2-m^2)R=0$，或 $R''+\frac1xR'+\big(1-\frac{m^2}{x^2}\big)R=0$ \\
贝塞尔级数 & $J_m(x)=\sum_{k\ge0}\frac{(-1)^k}{k!(k+m)!}\big(\frac x2\big)^{2k+m}$ \\
正交性 & $\int_0^aJ_m(k_{ml}r)J_m(k_{mn}r)r\dd r=\frac{a^2}2J_{m+1}^2(x_n^{(m)})\delta_{ln}$ \\
展开系数 & $A_n=\frac{2}{a^2J_{m+1}^2(x_n^{(m)})}\int_0^af(r)J_m(k_{mn}r)r\dd r$ \\
三类径向方程 & $\mu>0$ 贝塞尔；$\mu<0$ 修正贝塞尔；$\mu=0$ 欧拉 \\
\end{tabular}
\tabbot
\end{center}

\subhdr{6.2\quad 五条易错}

\begin{enumerate}[label=(\arabic*)]
  \item \textbf{把 $k$ 当成连续参数。}$k$ 必须让 $J_m(ka)=0$，所以 $k$ 是离散的一串 $x_n^{(m)}/a$。写答案时若 $k$ 还带“任意常数”的味道，就是没做完。
  \item \textbf{把 $k_n$ 写成只依赖 $n$。}零点依赖阶数：$x_1^{(0)}=2.405$，$x_1^{(1)}=3.832$。严格写 $k_{mn}$。
  \item \textbf{圆内保留 $N_m$。}$N_m$ 在 $r\to0$ 发散；圆内必须丢，环域才保留（环域不与 $r=0$ 相接）。
  \item \textbf{$\mu$ 的符号选错。}判据是“哪一组边界齐次”，见表 4-1；不是“贝塞尔函数好看就用 $\mu>0$”。
  \item \textbf{把波动的时间因子写成衰减。}波动是 $\cos(a kt)$、$\sin(a kt)$（两个初始条件，二阶时间导数）；输运是 $e^{-a^2k^2t}$（一个初始条件，一阶）。两者的空间谱相同，时间因子不同。
\end{enumerate}

\subhdr{6.3\quad 记号说明}

\begin{itemize}
  \item 课件 p3 写 $k=k_n=\dfrac{x_n^{(m)}}{a}$，p9 写 $\mu=\mu_n=\Big(\dfrac{x_n^{(m)}}{a}\Big)^2$：两处的 $n$ 都是“第 $n$ 个零点”，都隐含 $m$。本讲写 $k_{mn}$、$\mu_j$（用 $j$ 避免与 $n$ 冲突，见 \S5.2 的三重和）。
  \item 课件 p2 的径向方程中间有一处排版重叠（$(k^2r^2-m^2)RR=0$），实际是 $(k^2r^2-m^2)R=0$；本讲 \eqref{eq:radial} 按正确形式写。
  \item 课件 p10 把 $\mu=-\nu^2$ 的本征值写在底面条件上（$\nu_n=n\pi/h$），p9 把 $\mu>0$ 的本征值写在侧面条件上——两处的下标 $n$ 含义不同，读课件时要分开看。
\end{itemize}

%=====================================================================
\hdr{\S7}{课件页码对照表}
\phantomsection\label{sec:s7}
%=====================================================================

\begin{center}
\tabtop
\begin{tabular}{@{}L{0.08\linewidth}L{0.56\linewidth}L{0.30\linewidth}@{}}
\thead{课件页} & \thead{内容} & \thead{本讲位置} \\
\tabmidrule
p1 & 波动方程；$u=T(t)V$；$k^2$；Helmholtz；极坐标分离 & \S1.1--1.2 \\
p2 & $\lambda=m^2$；径向方程；$x=kr$；贝塞尔方程；$J_m,N_m$ & \S1.3 \\
p3 & $k=x_n^{(m)}/a$；$T$ 的方程；通解 & \S1.4--1.5 \\
p4 & 输运方程；$T'+a^2k^2T=0$；通解 & \S3.1 \\
p5 & 两次分离的总结；两组本征值 & \S3.1、\S6.1 \\
p6 & 作业 5、7 & \S9（P1、P2） \\
p7 & 柱坐标 Laplace；$u=R\Phi Z$；三个方程 & \S4.1 \\
p8 & $\lambda=m^2$；$\mu$ 的作用 & \S4.1--4.2 \\
p9 & $\mu>0$：贝塞尔；$\mu_j=(x_j^{(m)}/a)^2$；$Z=e^{\pm\sqrt\mu z}$ & \S4.3 \\
p10 & $\mu=-\nu^2$：修正贝塞尔；$\nu_n=n\pi/h$ & \S4.3 \\
p11 & $\mu=0$：欧拉方程；第二类侧面条件 & \S4.3 \\
p12 & 柱坐标通解；三种情形总结 & \S4.4 \\
p13 & Helmholtz 柱坐标；$Z''+\nu^2Z=0$；$\mu=k^2-\nu^2$ & \S5.1 \\
p14 & 三组本征值；$k=\sqrt{\mu+\nu^2}$ & \S5.2 \\
p15 & 三重级数；输运的指数衰减 & \S5.3 \\
\end{tabular}
\tabbot
\end{center}

%=====================================================================
\hdr{\S8}{三句话总结}
\phantomsection\label{sec:s8}
%=====================================================================

\begin{enumerate}[label=(\arabic*)]
  \item \textbf{贝塞尔函数不是新假设，是分离变量的必然结果。}圆域里出现时间，就出现 $k$；极坐标把 Helmholtz 拆开之后，径向方程 $x^2R''+xR'+(x^2-m^2)R=0$ 只能给出 $J_m$ 与 $N_m$。
  \item \textbf{谱由齐次边界条件给出。}圆内有界丢 $N_m$；$r=a$ 齐次给 $J_m(ka)=0$，即 $k=x_n^{(m)}/a$；底面齐次给 $n\pi/h$。不同边界给出不同的基，但正交性与展开公式是同一套。
  \item \textbf{柱坐标里 $\mu$ 的符号由“哪组边界齐次”决定：}侧面齐次 $\to$ 贝塞尔（$\mu>0$）；上下底齐次 $\to$ 修正贝塞尔（$\mu<0$）；侧面第二类且轴对称 $\to$ 欧拉（$\mu=0$）。
\end{enumerate}

%=====================================================================
\hdr{\S9}{独立练习（四题，含课件作业）}
\phantomsection\label{sec:practice}
%=====================================================================

P1、P2 是课件 p6 的作业 5 与作业 7（照录，不改编）；P3、P4 是本讲自编，把 \S2 的展开与 \S4 的判据各练一遍。详解在附录 A。

\practice{P1}{课件 p6 作业 5}{在半径 $\rho_0$ 的圆内解 $\nabla^2u=0$，边界条件分别为：(1) $u\big|_{r=\rho_0}=A\cos\varphi$；(2) $u\big|_{r=\rho_0}=A+B\sin\varphi$，其中 $A,B$ 为常数。分别求 $u$。}{圆内 Laplace 的 $r^m$ 项；边界值就是一个谐波时系数直接读出}{ans:P1}

\practice{P2}{课件 p6 作业 7}{内半径 $\rho_1$、外半径 $\rho_2$ 的环域内解稳态热传导 $\nabla^2u=0$；内圆温度 $u\big|_{r=\rho_1}=u_1\cos2\varphi$，外圆温度 $u\big|_{r=\rho_2}=u_2\sin\varphi$。求温度分布。}{环域保留 $r^{\pm m}$；两个不同阶的谐波各自独立定系数}{ans:P2}

\practice{P3}{自编，Fourier--Bessel 展开}{在 $0<r<a$ 上用 $J_0(x_n^{(0)}r/a)$ 展开 $f(r)=1-r^2/a^2$。(1) 写出系数公式；(2) 算出前三个系数（保留三位有效数字）；(3) 说明为什么这个级数在 $r=a$ 处给出 $0$，与 $f(a)=0$ 一致。}{Fourier--Bessel 系数公式；$\int_0^z tJ_0(t)\dd t=zJ_1(z)$}{ans:P3}

\practice{P4}{自编，$\mu$ 的判据}{下列四个柱坐标问题，分别应取 $\mu>0$、$\mu<0$ 还是 $\mu=0$？说明判据落在哪一组边界上。(a) 有限高圆柱，侧面 $u(a,\varphi,z)=0$，上下底给定非零温度；(b) 有限高圆柱，上下底 $u=0$，侧面给定温度；(c) 无限长圆柱（无 $z$ 边界），侧面 $u(a,\varphi)=f(\varphi)$；(d) 有限高圆柱，侧面绝热 $\partial_ru(a,\varphi,z)=0$，上下底给定温度，且温度与 $\varphi$ 无关。}{“哪组边界齐次”判据；$\mu=0$ 的轴对称欧拉分支}{ans:P4}

\clearpage
\appendix

%=====================================================================
\hdr{附录 A}{四道练习的详解}
\phantomsection\label{sec:app}
%=====================================================================

\answerstart{P1}{课件 p6 作业 5}{圆内 $r<\rho_0$，$\nabla^2u=0$；边界值分别为 $A\cos\varphi$ 与 $A+B\sin\varphi$。}

\textbf{第 1 步（圆内通解）。}由 \S8.2 的结论，圆内有界解是

\[
u(r,\varphi)=C_0+\sum_{m\ge1}r^m\big(a_m\cos m\varphi+b_m\sin m\varphi\big)
\]

（用 $a_m,b_m$ 以免与题给的常数 $A,B$ 混淆）。

\textbf{第 2 步（情形 (1)）。}边界值只含 $\cos\varphi$ 一项，故只有 $m=1$ 的余弦系数非零：

\[
a_1\rho_0=A,\qquad C_0=0,\qquad \text{其余为零},
\]

所以

\[
u(r,\varphi)=A\,\frac{r}{\rho_0}\cos\varphi=A\,\frac{x}{\rho_0}.
\]

\textbf{核对。}$\Delta u=0$（$x$ 是调和函数）；$r=\rho_0$ 时 $u=A\cos\varphi$，对。

\textbf{第 3 步（情形 (2)）。}边界值含常数与 $\sin\varphi$ 两项：

\[
C_0=A,\qquad b_1\rho_0=B,\qquad \text{其余为零},
\]

所以

\[
u(r,\varphi)=A+B\,\frac{r}{\rho_0}\sin\varphi=A+B\,\frac{y}{\rho_0}.
\]

\textbf{核对。}常数与 $y$ 都调和；$r=\rho_0$ 时回到 $A+B\sin\varphi$，对。两问都说明一件事：\textbf{边界值只含有限个谐波时，解只含那几项}，不需要写无穷级数。

\answerstart{P2}{课件 p6 作业 7}{环域 $\rho_1<r<\rho_2$，$\nabla^2u=0$；$u(\rho_1)=u_1\cos2\varphi$，$u(\rho_2)=u_2\sin\varphi$。}

\textbf{第 1 步（通解的形式）。}环域保留正负幂次：

\[
u(r,\varphi)=C_0+D_0\ln r+\sum_{m\ge1}\big(A_mr^m+C_mr^{-m}\big)\cos m\varphi+\sum_{m\ge1}\big(B_mr^m+D_mr^{-m}\big)\sin m\varphi .
\]

\textbf{第 2 步（筛项）。}两条边界值里只有 $\cos2\varphi$ 与 $\sin\varphi$，所以只保留 $m=2$ 的余弦与 $m=1$ 的正弦：

\[
u(r,\varphi)=\big(A_2r^2+C_2r^{-2}\big)\cos2\varphi+\big(B_1r+B_1'r^{-1}\big)\sin\varphi .
\]

（$C_0=D_0=0$：边界值没有常数项；其它谐波系数为零：正交性。）

\textbf{第 3 步（定 $\cos2\varphi$ 的系数）。}由 $r=\rho_1$ 与 $r=\rho_2$：

\[
\begin{cases}
A_2\rho_1^2+C_2\rho_1^{-2}=u_1,\\[2pt]
A_2\rho_2^2+C_2\rho_2^{-2}=0 .
\end{cases}
\]

第二式给 $C_2=-A_2\rho_2^4$；代入第一式：$A_2\big(\rho_1^2-\rho_2^4\rho_1^{-2}\big)=u_1$，即

\[
A_2=\frac{u_1\rho_1^2}{\rho_1^4-\rho_2^4},\qquad C_2=\frac{u_1\rho_1^2\rho_2^4}{\rho_2^4-\rho_1^4}.
\]

\textbf{第 4 步（定 $\sin\varphi$ 的系数）。}

\[
\begin{cases}
B_1\rho_1+B_1'\rho_1^{-1}=0,\\[2pt]
B_1\rho_2+B_1'\rho_2^{-1}=u_2,
\end{cases}
\quad\Longrightarrow\quad
B_1=\frac{u_2\rho_2}{\rho_2^2-\rho_1^2},\qquad B_1'=-\frac{u_2\rho_1^2\rho_2}{\rho_2^2-\rho_1^2}.
\]

\textbf{第 5 步（写出答案并核对）。}

\[
u=A_2\Big(r^2-\frac{\rho_2^4}{r^2}\Big)\cos2\varphi
+B_1\Big(r-\frac{\rho_1^2}{r}\Big)\sin\varphi ,
\]

其中 $A_2,B_1$ 如上。核对：$r=\rho_1$ 时第二项为零、第一项为 $A_2(\rho_1^2-\rho_2^4/\rho_1^2)\cos2\varphi=u_1\cos2\varphi$；$r=\rho_2$ 时第一项为零、第二项为 $B_1(\rho_2-\rho_1^2/\rho_2)\sin\varphi=u_2\sin\varphi$。两处都对。

\aside{本题最容易错的地方是把两个不同阶的谐波混在一个方程组里。它们的 $\varphi$ 依赖不同（$\cos2\varphi$ 与 $\sin\varphi$），由正交性各自独立，必须分开解两个 $2\times2$ 方程组。}

\answerstart{P3}{自编，Fourier--Bessel 展开}{$0<r<a$，$f(r)=1-r^2/a^2$，用 $J_0(x_n^{(0)}r/a)$ 展开。}

\textbf{第 1 步（系数公式）。}由 \eqref{eq:fb}，取 $m=0$、$k_{0n}=x_n^{(0)}/a$：

\[
A_n=\frac{2}{a^2J_1^2(x_n^{(0)})}\int_0^a\Big(1-\frac{r^2}{a^2}\Big)J_0\Big(\frac{x_n^{(0)}r}{a}\Big)r\dd r .
\]

\textbf{第 2 步（算积分）。}令 $t=x_n^{(0)}r/a$，则 $r=a t/x_n$，$r\dd r=(a/x_n)^2t\dd t$：

\[
\int_0^a\Big(1-\frac{r^2}{a^2}\Big)J_0\Big(\frac{x_nr}{a}\Big)r\dd r
=\Big(\frac{a}{x_n}\Big)^2\int_0^{x_n}\Big(1-\frac{t^2}{x_n^2}\Big)J_0(t)\,t\dd t .
\]

用 $\int_0^ztJ_0(t)\dd t=zJ_1(z)$、$\int_0^zt^2J_1(t)\dd t=z^2J_2(z)$ 与分部积分（$t^3J_0=\frac{\dd}{\dd t}[t^3J_1]-2t^2J_1$），得

\[
\int_0^{x_n}\Big(1-\frac{t^2}{x_n^2}\Big)J_0(t)t\dd t
=2J_2(x_n)
=\frac{4J_1(x_n)}{x_n}\quad(\text{用到 }J_0(x_n)=0,\ J_2=\tfrac{2}{x}J_1-J_0).
\]

\textbf{第 3 步（系数）。}把上式代回：

\[
A_n=\frac{2}{a^2J_1^2(x_n)}\cdot\Big(\frac{a}{x_n}\Big)^2\cdot\frac{4J_1(x_n)}{x_n}
=\frac{8}{x_n^3J_1(x_n)} .
\]

\textbf{第 4 步（前三个数值）。}用 $x_1=2.4048$、$J_1(x_1)=0.5192$；$x_2=5.5201$、$J_1(x_2)=-0.3403$；$x_3=8.6537$、$J_1(x_3)=0.2715$：

\begin{gather*}
A_1=\frac{8}{2.4048^3\times0.5192}=1.108,\qquad
A_2=\frac{8}{5.5201^3\times(-0.3403)}=-0.140,\\
A_3=\frac{8}{8.6537^3\times0.2715}=0.0455 .
\end{gather*}

系数按 $1/x_n^3$ 衰减，比 $f=1$ 那个例子快得多——因为 $1-r^2/a^2$ 在 $r=a$ 处等于零，没有端点跳跃。

\textbf{第 5 步（端点）。}$r=a$ 时每项 $J_0(x_n)=0$，级数为零；而 $f(a)=0$，两者一致，不存在跳跃。这与 \S2.5 的 $f=1$ 情形不同：那里 $f(a)=1$ 而级数给 $0$，端点必然跳跃，级数收敛到 $1/2$。

\answerstart{P4}{自编，$\mu$ 的判据}{四个柱坐标问题，判断 $\mu$ 的符号。}

判据只有一句：\textbf{哪一组边界是齐次的，就由哪一组定本征值；齐次条件要求对应的函数反复回到零，所以那一组必须是振荡型。}

\textbf{(a) 侧面 $u(a,\varphi,z)=0$（齐次），上下底给定（非齐次）。}侧面齐次 $\Rightarrow$ 径向函数要在 $r=a$ 为零 $\Rightarrow$ 贝塞尔，$\mu>0$；$Z$ 用 $e^{\pm\sqrt\mu z}$，其系数由上下底的非齐次数据定。

\textbf{(b) 上下底 $u=0$（齐次），侧面给定（非齐次）。}底面齐次 $\Rightarrow Z$ 要在 $z=0,h$ 为零 $\Rightarrow Z=\sin(n\pi z/h)$，$\mu=-\nu^2<0$，$\nu_n=n\pi/h$；径向方程是修正贝塞尔，$R=C I_m(\nu r)+D K_m(\nu r)$，实心柱丢 $K_m$。侧面条件用来定 $C$。

\textbf{(c) 无限长圆柱，无 $z$ 边界。}$z$ 方向没有边界条件，$Z$ 不能振荡也不能被某个 $z$ 条件定；取 $\mu=0$ 不损失一般性（$Z=A+Bz$，若问题与 $z$ 无关则 $Z$ 为常数）。这退化为 \S8.2 的二维极坐标问题：$\mu=0$，径向是欧拉方程。若题目明确 $u$ 与 $z$ 无关，就直接按二维做。

\textbf{(d) 侧面绝热（第二类齐次），上下底给定，轴对称。}$\varphi$ 无关 $\Rightarrow m=0$；侧面的第二类齐次条件 $\partial_ru(a)=0$ 对 $R=C_1+C_2\ln r$ 给 $C_2=0$，$R=C_1$ 自动满足；$Z=A+Bz$。这是 \S4.3 的 $\mu=0$ 分支。答案是一维线性温度 $u=A+Bz$，与半径无关——绝热侧面 + 轴对称的稳态，径向确实没有热流。

\aside{这道题的价值在于把“$\mu$ 的符号”从记公式变成判据：先问哪一组边界齐次，再问那一组要求的函数能不能反复回零，答案就唯一了。考试里若题目没有明说哪组齐次，就把所有的边界条件读一遍——非齐次的那一组是“数据入口”，齐次的那一组才是“谱的来源”。}

\end{document}
